\documentclass{controle}
\usepackage{main}

\title{Évaluation n°10 : Probabilités}
\author{Seconde 3}
\date{26 Janvier 2026}

\begin{document}
\titleandname{}
\version
\begin{questions}
\question[2] \hfill 

\vspace*{0.1cm}
Soient $A$ et $B$ deux événements d'une expérience aléatoire d'univers $\Omega$. Donner la formule de calcul de la probabilité $P(A \cup B)$.
\begin{equation*}
P(A \cup B) = \dots
\end{equation*}

\vspace*{0.25cm}
\question[8] \hfill

\vspace*{0.1cm}
Il y a dans une urne opaque \num{6} boules:
\begin{itemize}
\item \num{1} de couleur rouge, numérotée $1$;
\item \num{2} de couleur verte, numérotées de $1$ à $2$;
\item et \num{3} de couleur bleue, numérotées de $1$ à $3$. 
\end{itemize}
On tire alors au hasard deux boules de cette urne, \textbf{successivement et avec remise}. On note l'une après l'autre la couleur des boules obtenues.
\begin{parts}
\part[1] Quel est l'univers $\Omega$ de cette expérience aléatoire ?
\part[2] Représenter la situation à l'aide d'un \textbf{arbre de dénombrement}. Vous pouvez utiliser les notations de votre choix.
\part[1] Combien y a-t-il d'issues dans l'univers $\Omega$ ?
\part[2] Quelle est la probabilité de tirer deux boules de la même couleur ?
\part[2] Quelle est la probabilité de tirer deux boules de couleurs différentes ?
\end{parts}
\end{questions}

\newpage

\setcounter{page}{1}
\titleandname{}

\version
\begin{questions}
\question[2] \hfill 

\vspace*{0.1cm}
Soit $A$ un événement d'une expérience aléatoire d'univers $\Omega$. Donner la formule de calcul de la probabilité $P(\overbar{A})$.
\begin{equation*}
P(\overbar{A}) = \dots
\end{equation*}

\vspace*{0.25cm}
\question[8] \hfill

\vspace*{0.1cm}
Il y a dans une urne opaque \num{6} boules:
\begin{itemize}
\item \num{1} de couleur bleue, numérotée $1$;
\item \num{2} de couleur rouge, numérotées de $1$ à $2$;
\item et \num{3} de couleur verte, numérotées de $1$ à $3$. 
\end{itemize}
On tire alors au hasard deux boules de cette urne, \textbf{successivement et avec remise}. On note l'une après l'autre la couleur des boules obtenues.
\begin{parts}
\part[1] Quel est l'univers $\Omega$ de cette expérience aléatoire ?
\part[2] Représenter la situation à l'aide d'un \textbf{tableau à double entrée}. Vous pouvez utiliser les notations de votre choix.
\part[1] Combien y a-t-il d'issues dans l'univers $\Omega$ ?
\part[2] Quelle est la probabilité de tirer deux boules de la même couleur ?
\part[2] Quelle est la probabilité de tirer deux boules de couleurs différentes ?
\end{parts}
\end{questions}
\end{document}